\chapter{The Variable Vocabulary}\label{chap:6}

\begin{quote}
\emph{``The beginning of wisdom is to call things by their proper
name.''} --- Confucius
\end{quote}

Variables are the nouns of a CP-SAT model. Every decision the solver is
allowed to make is held by a variable, and every variable belongs to one
of five constructor families. The vocabulary is small by design: most of
the modeling work is recognizing which constructor matches a given
decision, and the recognition becomes reflex after a few examples.

CP-SAT has no continuous variables: there is no
\lstinline!new_real_var!. Quantities that are naturally
fractional (prices, distances, ratios) must be scaled to integers,
working in cents, millimeters, or some other unit small enough to
capture the precision the problem actually needs. The integer-only
restriction is part of the design.

\section{\texorpdfstring{Integer Variables --- \texttt{new\_int\_var}}{Integer Variables --- new\_int\_var}}\label{sec:6.1}

A bounded integer.

\begin{lstlisting}[language=Python]
x = model.new_int_var(lb=0, ub=100, name="x")  # closed [0, 100]
qty = model.new_int_var(0, 50, "qty")
delta = model.new_int_var(-10, 10, "delta")
\end{lstlisting}

The constructor takes a lower bound, an upper bound, and a name. Both
bounds are \emph{inclusive}: the variable's domain is the integer
interval \(\{lb, lb+1, \dots, ub\}\).

Tighter bounds tend to help. A bound is information the solver uses for
propagation, branching, and presolve. A variable declared as
\lstinline![0, 100]! is much cheaper to reason about than
the same variable declared as \([0, 10^9]\), even if the second bound is
technically valid. The looser bound forces the solver to consider
possibilities it never needed to.

Typical uses for integer variables include: quantities (\emph{how many
units to ship}, \emph{how many nurses on this shift}); times on a
discrete grid (\emph{the start hour of task A}, \emph{the minute at
which a meeting begins}); discrete spatial coordinates (\emph{the x- and
y-position of a piece on a grid}); counts and totals; indices into a
list (\emph{which warehouse serves customer C}, in
\(\{0, 1, \dots, n-1\}\)). For indices specifically, there is often a
competing encoding using one Boolean per choice (see \Cref{sec:6.2}).

\section{\texorpdfstring{Boolean Variables --- \texttt{new\_bool\_var}}{Boolean Variables --- new\_bool\_var}}\label{sec:6.2}

Booleans are the real workhorse of CP-SAT: the SAT core is built
around them. A Boolean is a 0/1 integer \emph{and} a logical literal at
the same time. This dual nature is what makes CP-SAT models concise.

\begin{lstlisting}[language=Python]
take = model.new_bool_var("take_item")
edge_used = model.new_bool_var("edge_a_b")
is_late = model.new_bool_var("is_late")
not_take = ~take  # negation — a free view, no new variable
\end{lstlisting}

The \lstinline!~! operator returns the negation as a
\emph{view}: no new variable is allocated, and the negation participates
in constraints exactly as the original Boolean would. This is a small
but useful efficiency: every Boolean in your model has a free
negation that can appear in any constraint position.

The dual nature is on display in expressions like:

\begin{lstlisting}[language=Python]
model.add(take + take2 + take3 <= 2)     # arithmetic: Bool as 0/1
model.add_bool_or([take, take2, take3])  # logic:      Bool as literal
\end{lstlisting}

The first treats the Booleans as integers and asserts a linear
inequality on their sum. The second treats them as logical literals and
asserts a disjunction. Both are legitimate, both are efficient, and the
same variable can appear in both kinds of constraint without conversion.

Typical uses for Booleans include: yes/no decisions (\emph{include
item}, \emph{use edge}, \emph{schedule task}, \emph{hire candidate});
assignments expressed as Bool-per-pair, the matrix idiom from the FLP
example of \Cref{sec:5.5}, with \(x_{i,j} \in \{0, 1\}\) indexed by pairs; state
indicators that other constraints react to (\emph{did the budget get
exceeded?}, \emph{is this nurse working a late shift?}); reified
relations, where a Boolean is constrained to be true exactly when some
other relation holds, used in conditional logic via
\lstinline!only_enforce_if!.

Booleans are typically a natural starting point for routing, scheduling,
and graph problems, where CP-SAT's SAT core is at its strongest. A
Bool-per-pair assignment encoding often works well because it exposes
many logical relations directly to the SAT layer. But it is not
universally better than a single integer ``which of \emph{k}'' variable:
integer-valued encodings can be superior when they enable stronger
global constraints or cleaner arithmetic. Treat the choice as a modeling
tradeoff and benchmark when it matters.

\section{\texorpdfstring{Constants --- \texttt{new\_constant}}{Constants --- new\_constant}}\label{sec:6.3}

A variable whose domain is a single value.

\begin{lstlisting}[language=Python]
zero = model.new_constant(0)
seven = model.new_constant(7)
\end{lstlisting}

We use this sparingly. The main case is \emph{interface uniformity}: a
function expects a variable, but for some inputs the value is fixed by
the data.

\begin{lstlisting}[language=Python]
def add_capacity(model, weight, capacity):
    model.add(weight <= capacity)

add_capacity(model, model.new_int_var(0, 100, "w"), 50)
add_capacity(model, model.new_constant(42), 50)
\end{lstlisting}

The second call passes a constant where the function expects a variable,
avoiding a separate code path for the fixed-value case. For most
expressions, CP-SAT also accepts a plain Python integer in place of a
variable, so \lstinline!new_constant! is genuinely needed
only when the API requires a variable object specifically.

Beyond uniformity, there is a modest debugging benefit:
\lstinline!new_constant! creates an actual
\lstinline!IntVar! (with an auto-generated name) that
appears in the exported model and the solver logs, where a bare integer
folded into a coefficient would not.

\section{\texorpdfstring{Interval Variables --- \texttt{new\_interval\_var}}{Interval Variables --- new\_interval\_var}}\label{sec:6.4}

An interval variable represents a span on a discrete axis: a
\emph{start}, a \emph{size} (duration), and an \emph{end}, with the
relation \lstinline!start + size = end! enforced
automatically. CP-SAT offers four flavors, increasing in expressiveness
and cost.

\begin{lstlisting}[language=Python]
# Fixed size, mandatory
fix = model.new_fixed_size_interval_var(start=s, size=5, name="fix")

# Variable size, mandatory
flex = model.new_interval_var(start=s, size=d, end=e, name="flex")

# Fixed size, optional (presence Boolean)
opt = model.new_optional_fixed_size_interval_var(
    s, 5, is_present=p, name="opt"
)

# Variable size, optional
opt_flex = model.new_optional_interval_var(
    s, d, e, is_present=p, name="opt_flex"
)
\end{lstlisting}

The \emph{mandatory} flavors require the interval to be present in any
feasible solution. The \emph{optional} flavors take an additional
Boolean \lstinline!is_present!. When
\lstinline!is_present = False!, the interval is treated
as \emph{absent} by every interval-aware constraint that touches it: it
does not occupy resources, does not interact with no-overlap
propagators, and effectively disappears from the scheduling problem for
that solution.

The cost rises with each step of expressiveness. A fixed-size mandatory
interval has the smallest representation; a variable-size optional
interval has the largest. The rule of thumb is to take the cheapest
flavor that expresses the problem.

Typical uses for intervals include: tasks on a timeline, with start,
duration, and end bundled into one object; optional tasks (\emph{this
maintenance might or might not happen this week}); rectangles in 2D
packing, where each rectangle is two intervals (one in x, one in y);
reservations on a shared resource (meetings, room bookings, machine
usage).

An interval is only interesting in combination with an interval-aware
constraint. By itself, an interval variable is a small bundle of three
integer variables with one arithmetic relation between them, no
special semantics. The reason intervals exist is that CP-SAT's
\emph{interval-aware constraints} (\lstinline!add_no_overlap!,
\lstinline!add_no_overlap_2d!,
\lstinline!add_cumulative!) accept intervals as input
and run dedicated scheduling propagators on the result.

If you only need a \lstinline![start, end]! span and never
feed it into one of these constraints, plain integer variables are
cheaper and clearer. The matching constraints are the next chapter's
subject; for now, treat intervals as the input shape to a scheduling
vocabulary you have not yet seen.

Intervals fit some shapes of scheduling better than others. They shine
when the time (or space) axis is \emph{somewhat continuous}: start
times anywhere in a wide horizon, sizes drawn from a range, optional
rectangles to be packed into a free 2D area. When the problem instead
has a small number of \emph{discrete slots} and tasks occupy known-sized
blocks, a Bool-per-(task, slot) encoding often models the same situation
more efficiently than intervals plus
\lstinline!add_no_overlap!. And when transitions between
tasks dominate the cost, sequence-dependent setup times, vehicle
routing, the problem is closer to routing than to scheduling, and
\lstinline!add_circuit! (next chapter) is the better
starting point.

\section{\texorpdfstring{Sparse Domains --- \texttt{new\_int\_var\_from\_domain}}{Sparse Domains --- new\_int\_var\_from\_domain}}\label{sec:6.5}

A variable whose domain is an explicit set of values rather than a
contiguous range.

\begin{lstlisting}[language=Python]
allowed = cp_model.Domain.from_values([2, 5, 8, 10, 20, 50, 90])
freq    = model.new_int_var_from_domain(allowed, "frequency_MHz")
\end{lstlisting}

This is the right tool when the restriction is \emph{part of the
problem}: allowed broadcast frequencies, standardized paper sizes,
container capacities, supported product configurations, the discrete set
of speeds at which a motor can run.

Using sparse domains as a \emph{performance} tool is also legitimate,
but it is a later move: CP-SAT uses \emph{lazy order encoding}
internally, so declaring
\lstinline!new_int_var(0, 100, ...)! does not eagerly
materialize one hundred Booleans, and a wide bound is not, in itself, a
wide cost. Pruning the domain to values you believe are reachable can
still help (sometimes substantially) but it can also hurt, and
worse, it can introduce subtle bugs when a value you excluded turns out
to be feasible after all. As a rule, start with the bounds the problem
gives you, and treat sparse-domain pruning as a tuning lever to
experiment with once a correctness baseline is in place.
